\chapter[24]{Heat Kernel for Evolving Metrics}
\label{ch:chow-iii-heat-kernel-evolving-metrics}
\dcref{ch24}
\source{chow-et-al-ricci-flow-part-iii:page-0286}

Let \(g(\tau)\), \(\tau\in[0,T]\), be a smooth family of complete metrics on
\(\mathcal M^n\), evolving by
\[
  \partial_\tau g_{ij}=2\mathcal R_{ij},
  \qquad \mathcal R=g^{ij}\mathcal R_{ij},
\]
and put
\[
  L_{x,\tau}=\partial_\tau-\Delta_{x,\tau}+Q,
  \qquad
  L^*_{y,v}=\partial_v+\Delta_{y,v}-Q+\mathcal R.
\]
Thus \(\partial_\tau d\mu_{g(\tau)}=\mathcal R\,d\mu_{g(\tau)}\).  The
backward Ricci-flow case is \(\mathcal R_{ij}=R_{ij}\) and \(Q=R\).

\section{Heat kernels and the local parametrix}

\begin{definition}[Heat kernel with respect to an evolving metric]
\label{def:chow-iii-evolving-heat-kernel}
\source{chow-et-al-ricci-flow-part-iii:page-0287}
\dcref{ch24:24.1}
\group{Evolving Heat Kernels}

A function \(H(x,\tau;y,v)\) on
\(\mathcal M\times\mathcal M\times\{0<\tau-v\le T\}\) is a
\emph{fundamental solution} for (L) if it has the standard space-time
regularity,
\[
 L_{x,\tau}H=0,\qquad L^*_{y,v}H=0,
\]
and converges to \(\delta_y\) as \(\tau\searrow v\), equivalently to
\(\delta_x\) as \(v\nearrow\tau\).  The minimal positive fundamental solution
is called the heat kernel.
\end{definition}

\begin{theorem}[Existence of the heat kernel on closed \(\mathcal M^n\)]
\label{thm:chow-iii-closed-evolving-heat-kernel}
\source{chow-et-al-ricci-flow-part-iii:page-0288}
\uses{def:chow-iii-evolving-heat-kernel, prop:chow-iii-evolving-parametrix, lem:chow-iii-parametrix-convolution-derivatives}
\dcref{ch24:24.2}
\group{Evolving Heat Kernels}


If \(\mathcal M\) is closed, then \(L\) has a unique fundamental solution
\(H(x,\tau;y,v)\).  It is positive and smooth, hence is the heat kernel.
\end{theorem}

For \(r_\tau(x)=d_{g(\tau)}(x,y)\), let
\[
 \widetilde E=(4\pi(\tau-v))^{-n/2}
 \exp\!\left(-\frac{r_\tau^2}{4(\tau-v)}\right),
 \qquad
 H_N=\widetilde E\sum_{k=0}^N\psi_k(\tau-v)^k.
\]
In polar coordinates about \(y\), write
\(\alpha_\tau=\sqrt{\det g^S(\tau)}/r_\tau^{n-1}\).

\begin{remark}[Changing distance in the parametrix]
\label{rem:chow-iii-changing-distance-parametrix}
\source{chow-et-al-ricci-flow-part-iii:page-0289}
\dcref{ch24:24.3}
\group{Local Parametrix}

Although \(\partial_\tau(r_\tau^2)\) is not usually explicit, only its local
expansion near \(y\) is needed in the parametrix construction.
\end{remark}

\begin{lemma}[Recursive ODEs for \(\psi_k\)]
\label{lem:chow-iii-parametrix-recursive-odes}
\source{chow-et-al-ricci-flow-part-iii:page-0289,chow-et-al-ricci-flow-part-iii:page-0290}
\uses{rem:chow-iii-changing-distance-parametrix}
\dcref{ch24:24.4}
\group{Local Parametrix}


The equations
\[
 \partial_{r_\tau}\psi_0=
 \left(-\tfrac12\partial_{r_\tau}\log\alpha_\tau
       +\tfrac12\partial_\tau r_\tau\right)\psi_0,
 \qquad \psi_0(y,y,\tau,v)=1,
\]
and, for \(k\ge1\),
\[
 r_\tau\partial_{r_\tau}\psi_k+
 \left(\frac{r_\tau}{2}\partial_{r_\tau}\log\alpha_\tau+k
       -\frac{r_\tau}{2}\partial_\tau r_\tau\right)\psi_k
 =-L_{x,\tau}\psi_{k-1}
\]
have unique finite solutions along radial geodesics.  They satisfy
\[
 L_{x,\tau}H_N=(\tau-v)^N\widetilde E,L_{x,\tau}\psi_N.
\]
\end{lemma}
\begin{proof}
\sketch Multiplication by the integrating factor
\[
 r_\tau^k\alpha_\tau^{1/2}
 \exp\!\left(-\frac12\int_0^{r_\tau}\partial_\tau r_\tau(\gamma(s))\,ds\right)
\]
reduces each equation to an integral along the radial geodesic.  Substitution
of the resulting recursion into \(L_{x,\tau}H_N\) cancels every power below
\((\tau-v)^N\).
\end{proof}

\begin{remark}[Independence of the truncation order]
\label{rem:chow-iii-parametrix-coefficients-independent}
\source{chow-et-al-ricci-flow-part-iii:page-0290}
\uses{lem:chow-iii-parametrix-recursive-odes}
\dcref{ch24:24.5}
\group{Local Parametrix}


The coefficient \(\psi_k\) is independent of the truncation order \(N\ge k\).
\end{remark}

\begin{remark}[Diagonal recursion]
\label{rem:chow-iii-parametrix-diagonal-recursion}
\source{chow-et-al-ricci-flow-part-iii:page-0291}
\uses{lem:chow-iii-parametrix-recursive-odes}
\dcref{ch24:24.6}
\group{Local Parametrix}


On the diagonal,
\[
 \psi_k(y,y,\tau,v)=-\frac1k
 (L_{x,\tau}\psi_{k-1})(x,y,\tau,v)\big|_{x=y}.
\]
\end{remark}

\begin{lemma}[Elementary bound for changing distances]
\label{lem:chow-iii-changing-distance-bound}
\source{chow-et-al-ricci-flow-part-iii:page-0291,chow-et-al-ricci-flow-part-iii:page-0292}
\dcref{ch24:24.7}
\group{Local Parametrix}

If \(\lvert\mathcal R_{ij}\rvert\le C\), then, away from the cut locus,
\[
 \left|\partial_\tau r_\tau(x)\right|\le C r_\tau(x).
\]
\end{lemma}
\begin{proof}
\sketch The first variation formula gives
\(
 \partial_\tau L_{g(\tau)}(\gamma)
 =\int_\gamma\mathcal R_\tau(\dot\gamma,\dot\gamma),ds
\)
for each minimizing geodesic.  Taking the infimum and using the tensor bound
proves the estimate.
\end{proof}

\begin{lemma}[Covariant derivatives of \(L_{x,\tau}H_N\)]
\label{lem:chow-iii-approximate-kernel-error-derivatives}
\source{chow-et-al-ricci-flow-part-iii:page-0292,chow-et-al-ricci-flow-part-iii:page-0293}
\uses{lem:chow-iii-parametrix-recursive-odes, lem:chow-iii-changing-distance-bound}
\dcref{ch24:24.8}
\group{Local Parametrix}


For \(k,\ell\ge0\), on the uniform injectivity neighborhood,
\[
 \left|\partial_\tau^\ell\nabla_x^k(L_{x,\tau}H_N)\right|
 \le C(\tau-v)^{N-n/2-k-2\ell}
 e^{-d_\tau^2(x,y)/(4(\tau-v))}.
\]
If \(N>n/2+k+2\ell\), the left side has the corresponding smooth
Gaussian-factor representation up to the diagonal.
\end{lemma}

Choose a cutoff \(\widetilde\eta\) supported where \(d_\tau(x,y)<2\rho\),
equal to one where \(d_\tau(x,y)\le\rho\), with
\(
 \rho=\frac14\min_{\tau\in[0,T]}\operatorname{inj}(g(\tau))
\), and set \(P_N=\widetilde\eta(d_\tau/\rho)H_N\).

\begin{definition}[Parametrix for \(L_{x,\tau}\)]
\label{def:chow-iii-evolving-parametrix}
\source{chow-et-al-ricci-flow-part-iii:page-0293}
\dcref{ch24:24.10}
\group{Global Parametrix}

A smooth function \(P\) is a parametrix if \(L_{x,\tau}P\) and
\(L^*_{y,v}P\) extend continuously to \(\tau=v\), and \(P\) converges to the
delta distribution in either time direction.
\end{definition}

\begin{proposition}[Existence of a parametrix for \(L_{x,\tau}\)]
\label{prop:chow-iii-evolving-parametrix}
\source{chow-et-al-ricci-flow-part-iii:page-0293}
\uses{def:chow-iii-evolving-parametrix, lem:chow-iii-approximate-kernel-error-derivatives}
\dcref{ch24:24.11}
\group{Global Parametrix}


If \(N>n/2\), then \(P_N\) is a parametrix for \(L_{x,\tau}\).
\end{proposition}

\begin{lemma}[Derivatives of \(L_{x,\tau}P_N\)]
\label{lem:chow-iii-cutoff-parametrix-error-derivatives}
\source{chow-et-al-ricci-flow-part-iii:page-0293}
\uses{prop:chow-iii-evolving-parametrix, lem:chow-iii-approximate-kernel-error-derivatives}
\dcref{ch24:24.12}
\group{Global Parametrix}


For smooth covariant tensors \(G_{k,\ell}\),
\[
 \partial_\tau^\ell\nabla_x^k(L_{x,\tau}P_N)
 =(\tau-v)^{N-n/2-k-2\ell}
 e^{-d_\tau^2(x,y)/(5(\tau-v))}G_{k,\ell}.
\]
In particular,
\[
 |L_{x,\tau}P_N|\le C_0(\tau-v)^{N-n/2}
 e^{-d_\tau^2(x,y)/(5(\tau-v))}.
\]
\end{lemma}

\section{The parametrix convolution series}

\begin{definition}[Space-time convolution]
\label{def:chow-iii-space-time-convolution}
\source{chow-et-al-ricci-flow-part-iii:page-0294}
\dcref{ch24:24.13}
\group{Convolution Series}

For suitable kernels \(F,G\), define
\[
 (F*G)(x,\tau;y,v)=\int_v^\tau\!\int_{\mathcal M}
 F(x,\tau;z,\sigma)G(z,\sigma;y,v)
 \,d\mu_{g(\sigma)}(z)\,d\sigma.
\]
\end{definition}

\begin{remark}[Static specialization of convolution]
\label{rem:chow-iii-static-convolution}
\source{chow-et-al-ricci-flow-part-iii:page-0294}
\dcref{ch24:24.14}
\group{Convolution Series}

When the metric is fixed and both kernels depend only on time differences,
Definition~\ref{def:chow-iii-space-time-convolution} is the usual heat-kernel
convolution.
\end{remark}

\begin{lemma}[Associativity of space-time convolution]
\label{lem:chow-iii-space-time-convolution-associative}
\source{chow-et-al-ricci-flow-part-iii:page-0294}
\uses{def:chow-iii-space-time-convolution}
\dcref{ch24:24.15}
\group{Convolution Series}


Whenever the relevant integrals are defined, space-time convolution is
associative.
\end{lemma}
\begin{proof}
\sketch Fubini's theorem identifies both bracketings with the integral over
\(v<\rho<\sigma<\tau\), after relabeling the two intermediate points.
\end{proof}

\begin{lemma}[First space derivatives of a convolution with \(P_N\)]
\label{lem:chow-iii-parametrix-convolution-first-derivative}
\source{chow-et-al-ricci-flow-part-iii:page-0295,chow-et-al-ricci-flow-part-iii:page-0296}
\uses{def:chow-iii-space-time-convolution, lem:chow-iii-cutoff-parametrix-error-derivatives}
\dcref{ch24:24.16}
\group{Convolution Series}


If \(G\) is continuous, then \(P_N*G\) is \(C^1\) in \(x\) and
\[
 \partial_{x^i}(P_N*G)=\int_v^\tau\!\int_{\mathcal M}
 (\partial_{x^i}P_N)(x,\tau;z,\sigma)G(z,\sigma;y,v)
 \,d\mu_{g(\sigma)}(z)\,d\sigma.
\]
\end{lemma}

\begin{lemma}[Second space derivatives of a convolution with \(P_N\)]
\label{lem:chow-iii-parametrix-convolution-second-derivative}
\source{chow-et-al-ricci-flow-part-iii:page-0296,chow-et-al-ricci-flow-part-iii:page-0297}
\uses{lem:chow-iii-parametrix-convolution-first-derivative}
\dcref{ch24:24.17}
\group{Convolution Series}


Under the hypotheses of Lemma~\ref{lem:chow-iii-parametrix-convolution-first-derivative},
in a geodesic coordinate neighborhood containing the half-injectivity ball,
\(P_N*G\) is \(C^2\) in \(x\), and its second derivatives are obtained by
differentiating \(P_N\) twice under the convolution integral.
\end{lemma}

\begin{lemma}[Time derivative of a convolution with \(P_N\)]
\label{lem:chow-iii-parametrix-convolution-time-derivative}
\source{chow-et-al-ricci-flow-part-iii:page-0297,chow-et-al-ricci-flow-part-iii:page-0298}
\uses{def:chow-iii-space-time-convolution, lem:chow-iii-cutoff-parametrix-error-derivatives}
\dcref{ch24:24.18}
\group{Convolution Series}


If \(G\) is continuous, then \(P_N*G\) is \(C^1\) in \(\tau\) and
\[
 \partial_\tau(P_N*G)=G+
 \int_v^\tau\!\int_{\mathcal M}(\partial_\tau P_N)G
 \,d\mu_{g(\sigma)}\,d\sigma.
\]
\end{lemma}

\begin{lemma}[Covariant derivatives of the convolution series]
\label{lem:chow-iii-parametrix-convolution-derivatives}
\source{chow-et-al-ricci-flow-part-iii:page-0298,chow-et-al-ricci-flow-part-iii:page-0299}
\uses{lem:chow-iii-space-time-convolution-associative, lem:chow-iii-parametrix-convolution-first-derivative, lem:chow-iii-parametrix-convolution-second-derivative, lem:chow-iii-parametrix-convolution-time-derivative}
\dcref{ch24:24.19}
\group{Convolution Series}


If \(m,\ell\ge0\) and \(N>n/2+2\ell+m\), then
\[
 \sum_{k=1}^\infty
 \partial_\tau^\ell\nabla_x^m\bigl((L_{x,\tau}P_N)^{*k}\bigr)
\]
converges absolutely and uniformly.  Thus its derivatives may be taken
termwise.  Consequently
\[
 H=P_N+P_N*\sum_{k=1}^\infty(L_{x,\tau}P_N)^{*k}
\]
is a smooth fundamental solution.
\end{lemma}

\section{Heat-kernel asymptotics}

\begin{theorem}[Heat kernel expansion for evolving metrics]
\label{thm:chow-iii-evolving-heat-kernel-expansion}
\source{chow-et-al-ricci-flow-part-iii:page-0299,chow-et-al-ricci-flow-part-iii:page-0300}
\uses{thm:chow-iii-closed-evolving-heat-kernel, lem:chow-iii-parametrix-convolution-derivatives}
\dcref{ch24:24.21}
\group{Heat Kernel Asymptotics}


On a closed manifold there are smooth \(u_k\), agreeing locally with
\(\psi_k\), such that
\[
 H=(4\pi(\tau-v))^{-n/2}e^{-d_\tau^2/(4(\tau-v))}
 \sum_{k=0}^N(\tau-v)^ku_k+w_N,
\]
where, uniformly in \(x,y\),
\[
 |\partial_\tau^\ell\nabla_x^kw_N|
 =O\bigl((\tau-v)^{N+1-n/2-k-2\ell}\bigr).
\]
\end{theorem}

\begin{lemma}[Expansion for \(\partial_\tau r_\tau\)]
\label{lem:chow-iii-changing-distance-expansion}
\source{chow-et-al-ricci-flow-part-iii:page-0301}
\uses{lem:chow-iii-changing-distance-bound}
\dcref{ch24:24.22}
\group{Heat Kernel Asymptotics}


In \(g(\tau)\)-normal coordinates \(x_\tau^i\) centered at \(y\),
\[
 \partial_\tau r_\tau(x)=\frac{\mathcal R_{ij}x_\tau^ix_\tau^j}{r_\tau}
 +\frac{(\nabla_k\mathcal R_{ij})x_\tau^kx_\tau^ix_\tau^j}{2r_\tau}
 +O(r_\tau^3),
\]
with tensors evaluated at \((y,\tau)\).
\end{lemma}

\begin{remark}[Integrated distance expansion]
\label{rem:chow-iii-integrated-distance-expansion}
\source{chow-et-al-ricci-flow-part-iii:page-0302}
\uses{lem:chow-iii-changing-distance-expansion}
\dcref{ch24:24.23}
\group{Heat Kernel Asymptotics}


Along the radial geodesic \(\gamma\), for \(0<s\le r_\tau(x)\),
\[
 \int_0^s\partial_\tau r_\tau(\gamma(\bar s))\,d\bar s
 =\frac{s^2}{2r_\tau(x)^2}\mathcal R_{ij}(y,\tau)x_\tau^ix_\tau^j
 +O(s^3).
\]
\end{remark}

\begin{lemma}[Asymptotic expansion for \(\psi_0\)]
\label{lem:chow-iii-psi-zero-expansion}
\source{chow-et-al-ricci-flow-part-iii:page-0302}
\uses{lem:chow-iii-parametrix-recursive-odes, lem:chow-iii-changing-distance-expansion}
\dcref{ch24:24.24}
\group{Heat Kernel Asymptotics}


In normal coordinates at (y),
\[
 \begin{aligned}
 \psi_0(x,y,\tau)={}&1+
 \left(\frac1{12}R_{ij}+\frac14\mathcal R_{ij}\right)x_\tau^ix_\tau^j\\
 &+\left(\frac1{24}\nabla_kR_{ij}
          +\frac1{12}\nabla_k\mathcal R_{ij}\right)
 x_\tau^kx_\tau^ix_\tau^j+O(r_\tau^4),
 \end{aligned}
\]
where the coefficients are evaluated at \((y,\tau)\).
\end{lemma}

\begin{lemma}[Asymptotic expansion for \(\psi_1\)]
\label{lem:chow-iii-psi-one-expansion}
\source{chow-et-al-ricci-flow-part-iii:page-0302,chow-et-al-ricci-flow-part-iii:page-0303,chow-et-al-ricci-flow-part-iii:page-0304,chow-et-al-ricci-flow-part-iii:page-0305,chow-et-al-ricci-flow-part-iii:page-0306}
\uses{lem:chow-iii-parametrix-recursive-odes, lem:chow-iii-psi-zero-expansion}
\dcref{ch24:24.25}
\group{Heat Kernel Asymptotics}


In normal coordinates at (y),
\[
 \begin{aligned}
 \psi_1(x,y,\tau)={}&
 \left(\frac16R+\frac12\mathcal R-Q\right)(y,\tau)\\
 &+\frac1{12}\left(\nabla_iR+2(\operatorname{div}\mathcal R)_i
 +\nabla_i\mathcal R-6\nabla_iQ\right)(y,\tau)x_\tau^i
 +O(r_\tau^2).
 \end{aligned}
\]
\end{lemma}

\section{Characterizing Ricci flow by heat-kernel asymptotics}

Let \(H=(4\pi\tau)^{-n/2}e^{-f}\) be the adjoint heat kernel based at
\((y,0)\), and set
\[
 \bar f=f-\int_{\mathcal M}fH\,d\mu_{g(\tau)},
 \qquad \Box=-\partial_\tau-\Delta_{g(\tau)}.
\]

\begin{claim}[Heat kernel asymptotics characterization of Ricci flow]
\label{clm:chow-iii-heat-kernel-characterizes-ricci-flow}
\source{chow-et-al-ricci-flow-part-iii:page-0307}
\uses{thm:chow-iii-evolving-heat-kernel-expansion, lem:chow-iii-psi-zero-expansion, lem:chow-iii-psi-one-expansion}
\dcref{ch24:24.26}
\group{Ricci Flow Characterization}


If \(\mathcal R_{ij}=R_{ij}\), so \(g(\tau)\) evolves by backward Ricci flow,
then
\[
 \Box\bar f(x,\tau)+\frac{\bar f(x,\tau)}{\tau}
 =O\bigl(d_{g(0)}(x,y)^2+\tau\bigr).
\]
In particular the left side is \(o(1)\) as
\(d_{g(0)}(x,y)^2+\tau\to0\).
\end{claim}

\begin{remark}[The cubic-over-time error]
\label{rem:chow-iii-cubic-over-time-error}
\source{chow-et-al-ricci-flow-part-iii:page-0310}
\dcref{ch24:24.27}
\group{Ricci Flow Characterization}

If \(r_\tau(x)^2+\tau\to0\) while \(r_\tau(x)^2/\tau\) remains bounded,
then \(r_\tau(x)^3/\tau\to0\).
\end{remark}

\begin{lemma}[Adjoint-kernel asymptotics under backward Ricci flow]
\label{lem:chow-iii-adjoint-parametrix-ricci-asymptotics}
\source{chow-et-al-ricci-flow-part-iii:page-0310}
\uses{lem:chow-iii-psi-zero-expansion, lem:chow-iii-psi-one-expansion, rem:chow-iii-cubic-over-time-error}
\dcref{ch24:24.28}
\group{Ricci Flow Characterization}


If \(\mathcal R_{ij}=R_{ij}\) and \(Q=R\), then for the approximate adjoint
heat kernel the quantity \(W\) defined from its logarithmic density satisfies
\[
 W=-\frac{n}{2\tau}
 +O\!\left(\frac{r_\tau(x)^3}{\tau}\right)
 +O\bigl(\tau+r_\tau(x)^2\bigr).
\]
\end{lemma}

\section{Dirichlet and noncompact heat kernels}

\begin{theorem}[Existence and uniqueness of Dirichlet heat kernel]
\label{thm:chow-iii-static-dirichlet-heat-kernel}
\source{chow-et-al-ricci-flow-part-iii:page-0311}
\dcref{ch24:24.30}
\group{Dirichlet Heat Kernels}

Let \((\mathcal M^n,g)\) be compact with nonempty smooth boundary.  There is a
unique nonnegative Dirichlet heat kernel \(H_D\), smooth in the interior,
vanishing on the boundary, and converging to \(\delta_y\) as \(t\searrow0\).
It has the eigenfunction expansion
\[
 H_D(x,y,t)=\sum_{k=1}^\infty e^{-\lambda_kt}\varphi_k(x)\varphi_k(y).
\]
\end{theorem}

\begin{lemma}[Dirichlet problem for the heat equation]
\label{lem:chow-iii-evolving-dirichlet-problem}
\source{chow-et-al-ricci-flow-part-iii:page-0312,chow-et-al-ricci-flow-part-iii:page-0313,chow-et-al-ricci-flow-part-iii:page-0314,chow-et-al-ricci-flow-part-iii:page-0315,chow-et-al-ricci-flow-part-iii:page-0316,chow-et-al-ricci-flow-part-iii:page-0317,chow-et-al-ricci-flow-part-iii:page-0318,chow-et-al-ricci-flow-part-iii:page-0319,chow-et-al-ricci-flow-part-iii:page-0320,chow-et-al-ricci-flow-part-iii:page-0321}
\uses{thm:chow-iii-closed-evolving-heat-kernel, lem:chow-iii-boundary-jump-relation, lem:chow-iii-normal-heat-kernel-derivative}
\dcref{ch24:24.31}
\group{Dirichlet Heat Kernels}


Let \(\mathcal M\) be compact with boundary and let \(b\) be continuous on
\(\partial\mathcal M\times(v,T]\), tending pointwise to zero as
\(\tau\searrow v\).  There is a unique continuous solution \(u\), smooth in
the interior, of
\[
 L_{x,\tau}u=0,\qquad u(\cdot,v)=0,
 \qquad u|_{\partial\mathcal M}=b.
\]
\end{lemma}
\begin{proof}
\sketch Extend the metric and potential across the double and let \(\widetilde H\) be
the resulting closed-manifold heat kernel.  Represent solutions by the normal
single-layer potential
\[
 u_\psi(x,\tau)=-\int_v^\tau\!\int_{\partial\mathcal M}
 \partial_{\nu_{z,\sigma}}\widetilde H(x,\tau;z,\sigma)
 \psi(z,\sigma)\,d\mu_{g(\sigma)}(z)\,d\sigma.
\]
The jump relation reduces the boundary condition to a Volterra integral
equation.  Picard iteration converges uniformly: the normal derivative has a
Gaussian estimate with integrable time singularity, and iterated kernels are
controlled by beta-function convolution.  The limit supplies \(u\), while the
maximum principle gives uniqueness.
\end{proof}

\begin{theorem}[Existence of Dirichlet heat kernel]
\label{thm:chow-iii-evolving-dirichlet-heat-kernel}
\source{chow-et-al-ricci-flow-part-iii:page-0313}
\uses{lem:chow-iii-evolving-dirichlet-problem, thm:chow-iii-closed-evolving-heat-kernel}
\dcref{ch24:24.32}
\group{Dirichlet Heat Kernels}


Let \(\widetilde H\) be the heat kernel on a closed extension of
\(\mathcal M\), and let \(f_{y,v}\) solve the homogeneous equation with zero
initial data and boundary value \(-\widetilde H(\cdot,\cdot;y,v)\).  Then
\[
 H(x,\tau;y,v)=\widetilde H(x,\tau;y,v)+f_{y,v}(x,\tau)
\]
is the fundamental solution satisfying the Dirichlet boundary condition.
\end{theorem}

\begin{lemma}[Jump relation]
\label{lem:chow-iii-boundary-jump-relation}
\source{chow-et-al-ricci-flow-part-iii:page-0313,chow-et-al-ricci-flow-part-iii:page-0314}
\uses{thm:chow-iii-closed-evolving-heat-kernel}
\dcref{ch24:24.33}
\group{Boundary Layer Potentials}


For \(x_0\in\partial\mathcal M\), the layer potential satisfies the
nontangential boundary limit
\[
 \lim_{x\to x_0}u_\psi(x,\tau)
 =\frac12\psi(x_0,\tau)-\int_0^\tau\!\int_{\partial\mathcal M}
 \partial_{\nu_{z,\sigma}}\widetilde H(x_0,\tau;z,\sigma)
 \psi(z,\sigma)\,d\mu_{g(\sigma)}(z)\,d\sigma.
\]
\end{lemma}

\begin{lemma}[Normal derivative estimate on a hypersurface]
\label{lem:chow-iii-normal-heat-kernel-derivative}
\source{chow-et-al-ricci-flow-part-iii:page-0317}
\uses{thm:chow-iii-evolving-heat-kernel-expansion}
\dcref{ch24:24.35}
\group{Boundary Layer Potentials}


Let \(\mathcal N^{n-1}\subset\mathcal M\) be a compact smooth hypersurface.
For \(x,y\in\mathcal N\), \(0\le\sigma<\tau\le T\), and a unit normal
\(\nu_{y,\sigma}\),
\[
 \left|\partial_{\nu_{y,\sigma}}H(x,\tau;y,\sigma)\right|
 \le C(\tau-\sigma)^{-n/2}
 \left(\frac{d_\tau^2(x,y)}{\tau-\sigma}+C\right)
 e^{-d_\tau^2(x,y)/(5(\tau-\sigma))}.
\]
It is consequently bounded by
\(C(\tau-\sigma)^{-n/2}e^{-d_\tau^2/(6(\tau-\sigma))}\).
\end{lemma}

\begin{remark}[Integrability of the first boundary kernel]
\label{rem:chow-iii-first-boundary-kernel-integrable}
\source{chow-et-al-ricci-flow-part-iii:page-0318}
\uses{lem:chow-iii-normal-heat-kernel-derivative}
\dcref{ch24:24.37}
\group{Boundary Layer Potentials}


For \(\alpha\in(1/2,1)\), the first iterated boundary kernel \(M_1\) obeys
\[
 \int_0^\tau\!\int_{\partial\mathcal M}|M_1|
 \,d\mu_{g(\sigma)}\,d\sigma\le C\tau^{1-\alpha},
\]
and hence its first Picard increment has the same bound.
\end{remark}

\begin{remark}[Gaussian normal derivative bounds]
\label{rem:chow-iii-normal-derivative-gaussian-bounds}
\source{chow-et-al-ricci-flow-part-iii:page-0318}
\uses{lem:chow-iii-normal-heat-kernel-derivative}
\dcref{ch24:24.38}
\group{Boundary Layer Potentials}


The two Gaussian estimates used for \(M_1\) follow directly from
Lemma~\ref{lem:chow-iii-normal-heat-kernel-derivative}.
\end{remark}

\begin{theorem}[Existence and uniqueness on a noncompact manifold, fixed metric]
\label{thm:chow-iii-noncompact-static-heat-kernel}
\source{chow-et-al-ricci-flow-part-iii:page-0322}
\dcref{ch24:24.39}
\group{Noncompact Heat Kernels}

Every, not necessarily complete, noncompact Riemannian manifold has a unique
minimal positive fundamental solution \(H_{\mathcal M}(x,y,t)\) of the heat
equation.  It is smooth and symmetric in \(x,y\).
\end{theorem}

\begin{theorem}[Existence and uniqueness on a noncompact manifold, evolving metric]
\label{thm:chow-iii-noncompact-evolving-heat-kernel}
\source{chow-et-al-ricci-flow-part-iii:page-0323,chow-et-al-ricci-flow-part-iii:page-0324}
\uses{thm:chow-iii-evolving-dirichlet-heat-kernel, lem:chow-iii-dirichlet-kernel-mass-bound}
\dcref{ch24:24.40}
\group{Noncompact Heat Kernels}


Let \(\mathcal M\) be noncompact, let \(g(\tau)\) be a smooth family of
metrics, and suppose \(Q\) is uniformly bounded.  Then \(L_{x,\tau}\) has a
unique smooth minimal positive fundamental solution
\(H_{\mathcal M}(x,\tau;y,v)\).
\end{theorem}
\begin{proof}
\sketch For an exhaustion \(\Omega_i\Subset\Omega_{i+1}\), the maximum principle gives
\(0<H_{\Omega_i}\le H_{\Omega_{i+1}}\).  The mass bound below and a local
parabolic mean-value inequality make the monotone limit finite on compact
subsets; local derivative estimates make it smooth.  It has the delta initial
condition and is dominated by every positive fundamental solution, proving
minimality and uniqueness.
\end{proof}

\begin{lemma}[Uniform mass bound for Dirichlet heat kernels]
\label{lem:chow-iii-dirichlet-kernel-mass-bound}
\source{chow-et-al-ricci-flow-part-iii:page-0323,chow-et-al-ricci-flow-part-iii:page-0324}
\uses{thm:chow-iii-evolving-dirichlet-heat-kernel}
\dcref{ch24:24.41}
\group{Noncompact Heat Kernels}


For \(y\in\Omega_i\), the Dirichlet heat kernel satisfies, uniformly in the
exhaustion index,
\[
 \int_{\Omega_i}H_{\Omega_i}(x,\tau;y,v)
 \,d\mu_{g(\tau)}(x)\le C(T)<\infty.
\]
\end{lemma}
\begin{proof}
\sketch Differentiate the integral, use the heat equation and the nonpositive outward
normal derivative on \(\partial\Omega_i\), and bound the potential and volume
evolution terms.  Gronwall's inequality and the limiting mass \(1\) at
\(\tau=v\) give a bound depending only on \(T\) and the uniform coefficient
bounds.
\end{proof}

\begin{claim}[Exercise: prove the jump relation]
\label{clm:chow-iii-boundary-jump-exercise}
\dcref{ch24:24.34}
\group{Boundary Layer Potentials}
\source{chow-et-al-ricci-flow-part-iii:page-0314}
\uses{lem:chow-iii-boundary-jump-relation}
Prove Lemma~\ref{lem:chow-iii-boundary-jump-relation} by isolating the
near-diagonal singularity of the normal derivative.
\end{claim}

\begin{claim}[Exercise: prove the error derivative estimate]
\label{clm:chow-iii-error-derivative-exercise}
\dcref{ch24:24.9}
\group{Local Parametrix}
\source{chow-et-al-ricci-flow-part-iii:page-0293}
\uses{lem:chow-iii-approximate-kernel-error-derivatives}
Prove Lemma~\ref{lem:chow-iii-approximate-kernel-error-derivatives} by
differentiating the Gaussian error representation.
\end{claim}

\begin{claim}[Exercise: complete the fundamental-solution construction]
\label{clm:chow-iii-fundamental-solution-exercise}
\dcref{ch24:24.20}
\group{Convolution Series}
\source{chow-et-al-ricci-flow-part-iii:page-0299}
\uses{lem:chow-iii-parametrix-convolution-derivatives}
Complete the details proving that the convolution series is a fundamental
solution of \(L_{x,\tau}u=0\) for the evolving metric.
\end{claim}

\begin{claim}[Exercise: prove the normal derivative estimate]
\label{clm:chow-iii-normal-derivative-exercise}
\dcref{ch24:24.36}
\group{Boundary Layer Potentials}
\source{chow-et-al-ricci-flow-part-iii:page-0317,chow-et-al-ricci-flow-part-iii:page-0318}
\uses{lem:chow-iii-normal-heat-kernel-derivative}
Prove Lemma~\ref{lem:chow-iii-normal-heat-kernel-derivative} using the
quadratic normal variation of squared distance along a hypersurface and the
time variation of its unit normal.
\end{claim}

\begin{claim}[Problem: prove the Ricci-flow characterization]
\label{clm:chow-iii-prove-ricci-flow-characterization}
\dcref{ch24:24.29}
\group{Ricci Flow Characterization}
\source{chow-et-al-ricci-flow-part-iii:page-0310}
\uses{clm:chow-iii-heat-kernel-characterizes-ricci-flow,lem:chow-iii-adjoint-parametrix-ricci-asymptotics}
Prove Claim~\ref{clm:chow-iii-heat-kernel-characterizes-ricci-flow} by
passing from the parametrix expansion to the heat kernel and normalizing (f).
\end{claim}
